\documentclass[../../main.tex]{subfiles}
\begin{document}

{\bf [解]}
\paragraph{(1)}

\begin{figure}[htb]
  \centering
  \begin{tikzpicture}[scale=3.6, >=stealth, every node/.style={font=\small}]

    % --- 1. 定数・幾何パラメータの定義 ---
    \def\R{1.0}
    \def\xVal{0.38}
    \pgfmathsetmacro{\yVal}{1.0 - \xVal}
    \pgfmathsetmacro{\rB}{\xVal / cos(30)}

    % 原点 O
    \coordinate (O) at (0, 0);

    % --- 2. 底面正六角形の頂点 B_k と側面の頂点 A_k の計算 ---
    \foreach \k in {1,...,6} {
      \pgfmathsetmacro{\angH}{90 + (\k - 1)*60}
      \pgfmathsetmacro{\angB}{60 + (\k - 1)*60}
      \coordinate (A\k) at (\angH:\R);
      \coordinate (B\k) at (\angB:\rB);
      \coordinate (H\k) at (\angH:\xVal);
    }

    % --- 3. 外接円 (半径 1) ---
    \draw[gray!60!black] (O) circle (\R);

    % --- 4. 底面正六角形 (B_1 ... B_6) ---
    \fill[teal!15] (B1) -- (B2) -- (B3) -- (B4) -- (B5) -- (B6) -- cycle;
    \draw[very thick, teal!85!black] (B1) -- (B2) -- (B3) -- (B4) -- (B5) -- (B6) -- cycle;

    % 中心 O から底面各頂点への線分
    \foreach \k in {1,...,6} {
      \draw[thin, teal!60!black] (O) -- (B\k);
    }

    % --- 5. 6つの側面二等辺三角形 (A_k B_k B_{k+1}) ---
    % カンマエラーを避けるためペアリストで指定
    \foreach \k/\nextk in {1/2, 2/3, 3/4, 4/5, 5/6, 6/1} {
      \fill[blue!8, opacity=0.7] (A\k) -- (B\k) -- (B\nextk) -- cycle;
      \draw[thick, blue!85!black] (B\k) -- (A\k) -- (B\nextk);
    }

    % --- 6. 注目する要素 (k = 1): 垂線 OA_1, 線分 OH_1, A_1H_1 ---
    \draw[thick, red!85!black, dashed] (O) -- (H1);
    \draw[thick, red!85!black] (H1) -- (A1);

    % 直角マーク (H1 における B1B2 ⊥ OA1)
    \draw[thin, gray!70]
    ($ (H1) + (-0.04, 0) $) -- ($ (H1) + (-0.04, 0.04) $) -- ($ (H1) + (0, 0.04) $);

    % 長さラベル x, y
    \node[left=2pt, red!85!black, font=\footnotesize] at ($ (O)!0.5!(H1) $) {$x$};
    \node[left=2pt, red!85!black, font=\footnotesize] at ($ (H1)!0.5!(A1) $) {$y$};

    % 底面の中心角 \pi/n の表示 (\angle B1-O-H1 = 30°)
    \draw[thin, purple!80!black] (0, 0.16) arc (90:60:0.16);
    \node[purple!90!black, font=\tiny, above right] at (75:0.17) {$\dfrac{\pi}{n}$};

    % --- 7. 主要点のプロットとラベル ---
    \fill (O) circle (0.5pt) node[below left=1pt] {$O$};
    \fill[red!85!black] (H1) circle (0.5pt) node[above right=1pt, red!85!black] {$H_1$};

    % 頂点 A_k
    \fill (A1) circle (0.6pt) node[above=2pt] {$A_1$};
    \fill (A2) circle (0.6pt) node[above left=1pt] {$A_2$};
    \fill (A3) circle (0.6pt) node[below left=1pt] {$A_3$};
    \fill (A4) circle (0.6pt) node[below=2pt] {$A_4$};
    \fill (A5) circle (0.6pt) node[below right=1pt] {$A_5$};
    \fill (A6) circle (0.6pt) node[above right=1pt] {$A_6$};

    % 底面頂点 B_k
    \fill[teal!85!black] (B1) circle (0.5pt) node[above=2pt, font=\scriptsize] {$B_1$};
    \fill[teal!85!black] (B2) circle (0.5pt) node[right=2pt, font=\scriptsize] {$B_2$};
    \fill[teal!85!black] (B3) circle (0.5pt) node[right=2pt, font=\scriptsize] {$B_3$};

    % 外接円の半径 1 の注記
    \node[gray!60!black, font=\scriptsize, right] at (10:1.02) {半径 $1$};

  \end{tikzpicture}
  \caption{正六角錐（$n=6$）の展開図と各幾何パラメータ（$OH_1=x,\ A_1H_1=y$）}
  \label{fig:pyramid_net_hexagon}
\end{figure}

図のように展開図での正$n$角錐の頂点を $A_1, A_2, \cdots, A_n$ とする．
又，底面の各点を $B_1 \cdots B_n$ とする．簡単のため$B_{n+1} = B_1$とする．
$A_k$ から $B_k B_{k+1}$ に下ろした垂足を $H_k$ とする．

\begin{figure}[htb]
  \centering
  \begin{tikzpicture}[scale=2.8, >=stealth, every node/.style={font=\small}]

    % --- 1. 定数・幾何パラメータの定義 ---
    % 手書き図の傾きに合わせて配置
    % OH_k 方向の角度: 32°
    % 半頂角: \pi/n = 30° (n = 6 の場合)
    % 高さ OH_k = x (図のスケールとして x = 2.0)
    \def\angH{32}
    \def\halfAng{30}
    \def\xVal{2.0}

    \coordinate (O) at (0, 0);

    % 垂足 H_k の座標
    \coordinate (H) at (\angH:\xVal);

    % 底面頂点 B_k, B_{k+1}
    % OB_k の偏角: 32° + 30° = 62°
    % OB_{k+1} の偏角: 32° - 30° = 2°
    \pgfmathsetmacro{\rB}{\xVal / cos(\halfAng)}
    \coordinate (Bk)    at ({\angH + \halfAng}:\rB);
    \coordinate (Bkone) at ({\angH - \halfAng}:\rB);

    % --- 2. 二等辺三角形 O B_k B_{k+1} の描画 ---
    \fill[teal!12] (O) -- (Bk) -- (Bkone) -- cycle;
    \draw[very thick, black] (O) -- (Bk) -- (Bkone) -- cycle;

    % --- 3. 垂線 OH_k (破線) ---
    \draw[thick, dashed, blue!85!black] (O) -- (H);

    % 直角マーク (H_k において OH_k ⊥ B_k B_{k+1})
    % H_k から O 方向への単位ベクトル、および B_k 方向への単位ベクトル
    \coordinate (u) at ($ 0.12*({cos(\angH + 180)}, {sin(\angH + 180)}) $);
    \coordinate (v) at ($ 0.12*({cos(\angH + 90)}, {sin(\angH + 90)}) $);
    \draw[thin, gray!70]
    ($ (H) + (u) $) -- ($ (H) + (u) + (v) $) -- ($ (H) + (v) $);

    % --- 4. 角度 \pi/n の表示 ---
    % 偏角 32° から 62° までの円弧
    \draw[purple!85!black, thick] (\angH:0.75) arc (\angH:{\angH + \halfAng}:0.75);
    \node[purple!90!black, font=\footnotesize, above left] at ({(\angH + 0.5*\halfAng)}:0.85) {$\dfrac{\pi}{n}$};

    % --- 5. 各種ラベルと寸法 ---
    % 高さ OH_k = x
    \node[blue!85!black, font=\footnotesize, below right=1pt] at ($ (O)!0.55!(H) $) {$x$};

    % 底辺 B_k B_{k+1} の寸法線（手書きの円弧状の引き出し線）
    \draw[<->, gray!80!black, thin]
    ($ (Bk) + (0.15, 0.08) $) to[bend left=35] ($ (Bkone) + (0.12, 0.08) $);
    \node[gray!90!black, font=\footnotesize, above right=1pt] at ($ (H) + (0.22, 0.12) $) {$2x\tan\dfrac{\pi}{n}$};

    % --- 6. 主要点のプロットとラベル ---
    \fill (O) circle (0.6pt) node[below left=2pt] {$O$};
    \fill (Bk) circle (0.6pt) node[above=2pt] {$B_k$};
    \fill (Bkone) circle (0.6pt) node[below right=2pt] {$B_{k+1}$};
    \fill[blue!85!black] (H) circle (0.6pt) node[above right=1pt, blue!85!black] {$H_k$};

  \end{tikzpicture}
  \caption{正$n$角錐の底面を構成する二等辺三角形 $\triangle O B_k B_{k+1}$}
  \label{fig:pyramid_base_triangle}
\end{figure}

$OH_1 = x, A_1 H_1 = y$ とおく ($0 < x < y$)と上図から
\begin{align}
  B_k B_{k+1} = 2x \tan \dfrac{\pi}{n} \label{eq:1}
\end{align}
と書ける．

次に，正$n$角錐の展開図が半径$1$の円に内接するとき，$OA_1$が円の半径と一致するので
\begin{align}
  |A_1O| = x + y = 1 \iff y = 1 - x \quad \Bigl(0 < x < \dfrac{1}{2}\Bigr) \label{eq:2}
\end{align}
である．

\begin{figure}[htb]
  \centering
  \begin{tikzpicture}[scale=2.4, >=stealth, every node/.style={font=\small}]

    % --- 1. 定数・幾何パラメータの定義 ---
    % x = OH_1, y = T H_1 (斜辺), h = OT (正n角錐の高さ)
    % 本文の増減表付近の代表値比率 (例: x = 0.4, y = 0.6 => h = sqrt(0.2) ≈ 0.447) を見やすく拡大
    \def\xVal{1.6}
    \def\hVal{2.2}
    \pgfmathsetmacro{\yVal}{sqrt(\xVal*\xVal + \hVal*\hVal)}

    % 座標の定義
    \coordinate (O)  at (0, 0);          % 底面の中心 (垂線の足)
    \coordinate (H1) at (\xVal, 0);      % 底面多角形の辺の中点
    \coordinate (T)  at (0, \hVal);      % 正n角錐の頂点

    % --- 2. 直角三角形の着色・描画 ---
    \fill[blue!8, opacity=0.7] (O) -- (H1) -- (T) -- cycle;
    \draw[very thick, black] (O) -- (H1) -- (T) -- cycle;

    % --- 3. 直角マーク (O における OT ⊥ OH_1) ---
    \draw[thin, gray!70]
    (0.18, 0) -- (0.18, 0.18) -- (0, 0.18);

    % --- 4. 辺の長さラベル ---
    % 底面中心から辺の中点まで: OH_1 = x
    \node[below=2pt, blue!90!black] at ($ (O)!0.5!(H1) $) {$x$};

    % 正n角錐の高さ: OT = h
    \node[left=2pt, red!85!black] at ($ (O)!0.5!(T) $) {$h$};

    % 側面の高さ (母線上の垂線): TH_1 = y
    \node[above right=1pt, teal!85!black] at ($ (H1)!0.5!(T) $) {$y$};

    % --- 6. 主要点のプロットと頂点ラベル ---
    \fill (O)  circle (0.6pt) node[below left=2pt] {$O$};
    \fill (H1) circle (0.6pt) node[below right=2pt] {$H_1$};
    \fill[red!85!black] (T) circle (0.7pt) node[above=2pt, red!85!black] {頂点 $T$};

  \end{tikzpicture}
  \caption{正$n$角錐の軸を含む断面（直角三角形 $\triangle TOH_1$）}
  \label{fig:pyramid_vertical_section}
\end{figure}

ここで正$n$角錐の高さを$h$として上図の三角形$OTH_1$に三平方の定理を用いて
\begin{align}
  h = \sqrt{y^2 - x^2} = \sqrt{(1-x)^2 - x^2} = \sqrt{1 - 2x} \label{eq:3}
\end{align}
を得る．ここで二つ目の等号に\cref{eq:2}を用いた．

次に正$n$角錐の底面積 $S$ は
\begin{align}
  S
  &= n \times (\triangle O B_1 B_2) \\
  &= n \times \dfrac{|B_{k}B_{k+1||OH_{k}|}}{2} \\
  &= n \times \dfrac{1}{2} \cdot x \cdot \Bigl(2x \tan\dfrac{\pi}{n}\Bigr) \\
  &= n x^2 \tan\dfrac{\pi}{n}
\end{align}
となる．従って\cref{eq:3}と合わせて，体積 $V$ は
\begin{align}
  V = \dfrac{1}{3} S \cdot h = \dfrac{1}{3} n \tan\dfrac{\pi}{n} \, x^2 \sqrt{1 - 2x} \label{eq:4}
\end{align}
である．$x$を動かした時の$V$の最大値を求めるため，$V$を$x$で微分すると
\begin{align}
  \dfrac{dV}{dx}
  &= \dfrac{1}{3} n \tan\dfrac{\pi}{n} \left(4x^3-10x^4\right) \\
  &= \dfrac{1}{3} n \tan\dfrac{\pi}{n} 2x^3\left(2-5x\right)
\end{align}
である．従って$V$の増減表は下表のようになる．

\begin{table}[htb]
  \centering
  \caption{$V$ の増減表}
  \begin{tabular}{c|c|c|c|c|c}
    $x$ & $0$ & $\dots$ & $\dfrac{2}{5}$ & $\dots$ & $\dfrac{1}{2}$ \\ \hline
    $V'$ & & $+$ & $0$ & $-$ & \\ \hline
    $V$ & & $\nearrow$ & 極大 & $\searrow$ &
  \end{tabular}
\end{table}

従って $x = \dfrac{2}{5}$ の時に $V$ も最大になり，その値は
\begin{align}
  V &= v_n \\
  &= \dfrac{1}{3} n \tan\dfrac{\pi}{n} \cdot \Bigl(\dfrac{2}{5}\Bigr)^2 \sqrt{\dfrac{1}{5}} \\
  &= \dfrac{4\sqrt{5}}{375} n \tan\dfrac{\pi}{n} \label{eq:4}
\end{align}
である．

\paragraph{(2)}

$t = \dfrac{\pi}{n}$ とおくと $n \to \infty$ の時に $t \to 0$ だから，求める極限値は
\begin{align}
  \lim_{n\to\infty}v_n
  &= \lim_{t\to 0}\dfrac{4\sqrt{5}}{375} \pi \dfrac{\tan t}{t} \\
  &= \dfrac{4\sqrt{5}}{375} \pi
\end{align}
である．

\end{document}
